\documentclass[11pt,a4paper]{article}
\usepackage{amsmath,amssymb}
\usepackage{listings}
\usepackage{hyperref}
\usepackage[margin=1in]{geometry}
\usepackage{fancyhdr}

\pagestyle{fancy}
\fancyhf{}
\fancyhead[L]{Module 14: GTK4 GUI Programming}
\fancyhead[R]{C Programming Course}
\fancyfoot[C]{\thepage}

\lstset{
  language=C,
  basicstyle=\ttfamily\small,
  keywordstyle=\bfseries,
  commentstyle=\itshape,
  numbers=left,
  numberstyle=\tiny,
  frame=single,
  breaklines=true,
  tabsize=4
}

\title{Module 14: GTK4 GUI Programming}
\author{C Programming Course}
\date{}

\begin{document}
\maketitle
\thispagestyle{fancy}

\section{Introduction to GTK4}
GTK4 is a cross-platform toolkit for creating graphical user interfaces. It
provides widgets, layout containers, and an event-driven programming model.

\section{Setting Up GTK4 Development}
On Debian/Ubuntu:
\begin{verbatim}
sudo apt-get install libgtk-4-dev
\end{verbatim}

Compile GTK4 programs with \texttt{pkg-config}:
\begin{verbatim}
gcc `pkg-config --cflags gtk4` -o app app.c \
    `pkg-config --libs gtk4`
\end{verbatim}

\section{Creating Windows and Widgets}
\begin{lstlisting}
#include <gtk/gtk.h>

static void activate(GtkApplication *app,
                     gpointer user_data) {
    GtkWidget *window = gtk_application_window_new(app);
    gtk_window_set_title(GTK_WINDOW(window), "Hello");
    gtk_window_set_default_size(
        GTK_WINDOW(window), 400, 300);

    GtkWidget *label = gtk_label_new("Hello, GTK4!");
    gtk_window_set_child(GTK_WINDOW(window), label);
    gtk_window_present(GTK_WINDOW(window));
}

int main(int argc, char **argv) {
    GtkApplication *app = gtk_application_new(
        "org.example.hello",
        G_APPLICATION_DEFAULT_FLAGS);
    g_signal_connect(app, "activate",
                     G_CALLBACK(activate), NULL);
    int status = g_application_run(
        G_APPLICATION(app), argc, argv);
    g_object_unref(app);
    return status;
}
\end{lstlisting}

\section{Event Handling}
GTK4 uses signals and callbacks. Connect a signal to a handler with
\texttt{g\_signal\_connect}:

\begin{lstlisting}
static void on_button_clicked(GtkButton *btn,
                              gpointer data) {
    g_print("Button clicked!\n");
}
/* In activate: */
/* g_signal_connect(button, "clicked",
       G_CALLBACK(on_button_clicked), NULL); */
\end{lstlisting}

\section{GTK4 Best Practices}
\begin{itemize}
  \item Use \texttt{GtkBox} or \texttt{GtkGrid} for layouts.
  \item Keep UI logic in the activate callback.
  \item Unreference objects when no longer needed.
  \item Use GLib main loop for asynchronous operations.
\end{itemize}

\end{document}
